\section*{9. 電池}
\begin{enumerate}
    \item 電池とはどのような装置か．次の語句をすべて用いて説明せよ．\\
    語句：化学エネルギー，電気エネルギー\\
    (\hspace{15cm})
    \item 次のア～エの水溶液を電池の水溶液として使用したとき，電流が流れるものをすべて選べ．\\
    ア：水酸化ナトリウム水溶液\hspace{1cm}イ：塩化ナトリウム水溶液\hspace{1cm}\\
    ウ：精製水\hspace{1cm}エ：エタノール\\[0.2cm]
    答(\hspace{3cm})
    \item うすい塩酸の入っているビーカーに銅板とマグネシウム板を入れて電池を作った．このときの+極と-極はそれぞれどちらの金属板になるか．また，その理由を答えろ．\\
    +極(\hspace{2cm})\hspace{1cm}-極(\hspace{2cm})\\
    理由(\hspace{15cm})
    \item ダニエル電池において，ビーカー中にセロハンを置く理由を二つ答えよ．
    \begin{itemize}
        \item[1] (\hspace{15cm})\vspace{0.5em}
        \item[2] (\hspace{15cm})
    \end{itemize}
    \item \begin{enumerate}
        \item 使うと電圧が下がり，元に戻らない電池を何というか．またその例を1つかけ．\\
        (\hspace{3cm})\hspace{1cm}例(\hspace{2cm})
        \item 使うと電圧が下がるが，外部から逆向きの電流をかけることで電圧が回復し，再び使用可能となる電池を何というか．またその例を1つかけ．\\
        (\hspace{3cm})\hspace{1cm}例(\hspace{2cm})
    \end{enumerate}
    \item 水の電気分解の逆の操作を行うことで電力を取り出す電池を何というか．また，その化学反応式を示せ．\\
    (\hspace{2cm})\hspace{1cm}化学反応式(\hspace{8cm})
\end{enumerate}
